\documentclass[10pt,a4paper]{amsart}
\usepackage[margin=19mm]{geometry}
\usepackage{amsmath,amssymb,mathtools}
\usepackage[scr=rsfs,cal=euler]{mathalfa}
\usepackage{xfrac}
\usepackage[unicode,hidelinks,bookmarks=false]{hyperref}
\newcommand{\ie}{i.e.\ }
\newcommand{\cf}{cf.\ }
\newcommand{\resp}{respectively\ }
\newcommand{\Subsection}{\subsection}
\providecommand{\nobreakdash}{\nobreak\hbox{-}}
\setlength{\emergencystretch}{2em}
\multlinegap0pt
\usepackage{bm}
\usepackage{bbm}
\usepackage{dsfont}
\usepackage{xspace}
\usepackage{cancel}
\usepackage{mathtools}
\usepackage{amsthm}
\makeatletter
\@ifundefined{theorem}{\newtheorem{theorem}{Theorem}}{}
\@ifundefined{lemma}{\newtheorem{lemma}[theorem]{Lemma}}{}
\makeatother
\title[A $\mathcal{CR}$-rotated $Q\_1$ nonconforming finite element ]{A $\mathcal{CR}$-rotated $Q\_1$ nonconforming finite element method for Stokes interface problems on local anisotropic fitted mixed meshes}
\author{Geng Chenchen, Hua Wang, Fengren Zou}
\date{Source: arXiv:2507.02741v1 (CC BY 4.0)}
\begin{document}
\maketitle

\noindent\emph{Source and licence.} A $\mathcal{CR}$-rotated $Q_1$ nonconforming finite element method for Stokes interface problems on local anisotropic fitted mixed meshes. Version arXiv:2507.02741v1 (CC BY 4.0), distributed under the Creative Commons Attribution 4.0 International licence, as recorded in the arXiv licence metadata for this exact version and shown on the abstract page https://arxiv.org/abs/2507.02741v1. Full rights notice: licenses/arxiv-2507.02741-CC-BY-4.0.txt.\par\smallskip

\noindent\textit{Editorial scope.} Counted unit: the Lemma whose source label is \texttt{le:sup1} in the article (source0.tex lines 911--916, source label \texttt{le:sup1}), delivered together with the article's own complete original proof of it (source0.tex lines 918--999, one \texttt{proof} environment of 82 lines). No mathematics is written, repaired or re-derived: statement and proof are the article's own text line for line. Required context retained uncounted: lem:basis-est (lines 717--723); le:div (lines 863--870). Every definition, display and label of those lines is retained unchanged. Omitted from this package and not redistributed: 1--716, 724--862, 871--910, 917--917, 1000--1362. Nothing inside a retained excerpt is omitted or altered: every retained body is delivered exactly as the article's lines are written, so no omission receipt of any kind accompanies this package. Citations: the retained excerpts carry no citation command at all, so no bibliography entry of the article is transcribed and no reference is redistributed. Reference adaptation: none was needed and none was made; every reference inside the delivered unit resolves inside it, so no unresolved reference or citation is produced. Printed slips kept literal: no printed word, symbol, hypothesis or number of the counted statement or of its proof was altered. Layout and class: the article's own class is \texttt{article} with packages including amsmath, amssymb, amsthm, amsfonts, graphicx, epstopdf, array, booktabs, tabularx, subfigure, caption, float, color, bm; the wrapper is the lane's pinned \texttt{amsart} editorial wrapper at 10pt on A4, which restates the theorem-like environments the retained text needs and adds the article's own macro definitions that the retained text needs (none), with extra wrapper packages (bm, bbm, dsfont, xspace, cancel, mathtools). The delivered numbering is the unit's own. Assets: no retained span contains a figure environment, a graphics-inclusion command, a PDF-page inclusion, a table or a TikZ picture; no image file is copied into this package, the package declares no asset dependency and no image file is redistributed with it. Licence: version arXiv:2507.02741v1 of this article is distributed under the Creative Commons Attribution 4.0 International licence, as recorded in the arXiv licence metadata for this exact version and shown on the abstract page https://arxiv.org/abs/2507.02741v1. A full rights notice is delivered at licenses/arxiv-2507.02741-CC-BY-4.0.txt. Characters: every retained excerpt is pure ASCII, so the delivered engine, XeLaTeX with its default Latin Modern text font, reports no missing character.
\begin{lemma}\label{lem:basis-est}
	Let $M$ be one of the subelement of the interface macro element, it holds that
	\begin{align}
		\|\partial_{\tilde{x}} \tilde{\phi}_{\Gamma,e}\|_{L^2(\tilde{K})} \lesssim |K|^{1/2},\\
		\|\partial_{\tilde{y}} \tilde{\phi}_{\Gamma,e}\|_{L^2(\tilde{K})} \lesssim |K|^{-1/2}.
	\end{align} 
\end{lemma}

\begin{lemma}\label{le:div}
	For interface elements $M \in \mathcal{T}_h^{\Gamma}$, there exists a positive constant $C$ independent of $h$, such that for all $p \in L^2_0(M)$, there exists $\bm{v} \in (H^1_0(M))^2$ satisfying:
	\begin{align}
		\nabla \cdot \bm{v} &= p \quad \mathrm{in}\, M, \\
		|\bm{v}|_{H^1(M)} &\leq C \|p\|_{L^2(M)},
	\end{align}
	where $C$ depends only on the maximum angle condition of the triangulation.
\end{lemma}

\begin{lemma}\label{le:sup1}
	There exists a constant $k_1 > 0$, independent of $h$, such that for any $q_{0,h} \in X_{0,h}$, it holds
	\begin{equation}
		k_1\| q_{0,h} \|_{L^2(\Omega)} \leq \sup_{\substack{\boldsymbol{v}_{1,h} \in \boldsymbol{U}_{0,h} \\ \boldsymbol{v}_{1,h} \neq \bm{0}}} \frac{b_h(\boldsymbol{v}_{1,h}, q_{0,h})}{\|\boldsymbol{v}_{1,h}\|_{\boldsymbol{U}_{h}}}.
	\end{equation}
\end{lemma}

\begin{proof}
	For any $q_{0,h} \in X_{0,h}$ and any macro-element $M \in \mathcal{T}_h^{\Gamma}$, since $\int_M q_{0,h}  dx = 0$, Lemma \ref{le:div} guarantees the existence of $\boldsymbol{v}_1 \in \bm{H}^1_0(M)$ satisfying:
	\begin{align*}
		\nabla \cdot \boldsymbol{v}_1 &= q_{0,h}\quad \text{in} ~M, \\
		|\boldsymbol{v}_1|_{H^1(M)} &\lesssim \|q_{0,h}\|_{L^2(M)}.
	\end{align*}
	
	Define the interpolation operator $\Pi^{(1)}_M: \bm{H}^1(M) \rightarrow \bm{U}_{h}(M)$ by
	\begin{equation}\label{eq:pi_m^1}
		\int_e \Pi^{(1)}_M \boldsymbol{v}  ds = \int_e \boldsymbol{v}  ds, \quad \forall e \in \mathcal{E}(M).
	\end{equation}
	This implies the representation:
	\[
	\Pi^{(1)}_M \boldsymbol{v} = \sum_{i=1}^{6} \left( \frac{1}{|e_i|} \int_{e_i} \boldsymbol{v}  ds \right) \boldsymbol{\phi}_i.
	\]
	For $\boldsymbol{v}_1 \in \bm{H}^1_0(M)$, we have specifically:
	\[
	\Pi^{(1)}_M \boldsymbol{v}_1 = \left( \frac{1}{|e_6|} \int_{e_6} \boldsymbol{v}_1  ds \right) \boldsymbol{\phi}_6.
	\]
	
	We first establish the stability on the triangular sub-element $T$:
	Since $\Pi^{(1)}_M \boldsymbol{v}|_T \in \bm{P}_1(T)$, its gradient is constant. By Green's formula:
	\[
	\int_T \partial_x (\Pi^{(1)}_M \boldsymbol{v})  dx = \int_T \partial_x \boldsymbol{v}  dx.
	\]
	Thus,
	\[
	\partial_x (\Pi^{(1)}_M \boldsymbol{v})|_T = \frac{1}{|T|} \int_T \partial_x \boldsymbol{v}  dx.
	\]
	The $L^2$-norm satisfies:
	\begin{align*}
		\|\partial_x (\Pi^{(1)}_M \boldsymbol{v})\|_{L^2(T)} 
		&= |T|^{1/2} \left| \frac{1}{|T|} \int_T \partial_x \boldsymbol{v}  dx \right| \\
		&\leq |T|^{-1/2} \left| \int_T \partial_x \boldsymbol{v}  dx \right| \\
		&\lesssim \|\partial_x \boldsymbol{v}\|_{L^2(T)}.
	\end{align*}
	Similarly, $\|\partial_y (\Pi^{(1)}_M \boldsymbol{v})\|_{L^2(T)} \lesssim \|\partial_y \boldsymbol{v}\|_{L^2(T)}$. Therefore,
	\begin{equation}\label{eq:pi_m^1err}
		|\Pi^{(1)}_M \boldsymbol{v}|_{H^1(T)} \lesssim |\boldsymbol{v}|_{H^1(T)}.
	\end{equation}
	
	For the quadrilateral sub-element $Q$, using scaling argument, we have
	\begin{equation*}
		|\Pi^{(1)}_M \boldsymbol{v}|_{H^1(Q)} \lesssim |\tilde{\Pi}^{(1)}_{\tilde{M}} \boldsymbol{\tilde{v}}|_{H^1(\tilde{Q})}.
	\end{equation*}
    and 
    \[
    \tilde{\Pi}^{(1)}_{\tilde{M}} \boldsymbol{\tilde{v}}_1 = \left( \frac{1}{|\tilde{e}_6|} \int_{\tilde{e}_6} \boldsymbol{\tilde{v}}_1  ds \right) \boldsymbol{\tilde{\phi}}_6.
    \]
    since $\partial_{\tilde{y}} (\tilde{\Pi}^{(1)}_{\tilde{M}} \boldsymbol{\tilde{v}}_1)$ is constant in $\tilde{Q}$, using the same argument as for the triangle subelement, we have 
    \[
    \|\partial_{\tilde{y}} (\tilde{\Pi}^{(1)}_{\tilde{M}} \boldsymbol{\tilde{v}}_1)\|_{L^2(Q)} \lesssim \|\partial_{\tilde{y}} \boldsymbol{\tilde{v}}_1\|_{L^2(\tilde{Q})}.
    \]
    
    Using lemma \ref{lem:basis-est}, we have
    \begin{align*}
    	|\partial_{\tilde{x}}\tilde{\phi}_6|_{H^1(\tilde{Q})}\lesssim |\tilde{Q}|^{1/2}.
    \end{align*}
    consequently,  let $\delta \boldsymbol{\tilde{v}}_1 = \boldsymbol{\tilde{v}}_1 - \frac{1}{|\tilde{Q}|}\int_{\tilde{Q}}\boldsymbol{\tilde{v}}_1 d\tilde{x}$, we derive
	\begin{align*}
		\|\partial_{\tilde{x}} (\Pi^{(1)}_{\tilde{M}} \boldsymbol{\tilde{v}}_1)\|_{L^2(\tilde{Q})} 
		&=\|\partial_{\tilde{x}} (\Pi^{(1)}_{\tilde{M}} \delta\boldsymbol{\tilde{v}}_1)\|_{L^2(\tilde{Q})} \\
		&= \left| \frac{1}{|\tilde{e}_6|} \int_{\tilde{e}_6} \delta\boldsymbol{\tilde{v}}_1  ds \right| \|\partial_{\tilde{x}} \boldsymbol{\tilde{\phi}}_6\|_{L^2(\tilde{Q})} \\
		&\lesssim \frac{|\tilde{Q}|^{1/2}}{|\tilde{e}_6|^{1/2}}\|\delta\boldsymbol{\tilde{v}}_1\|_{L^2(\tilde{e}_6)}\\
		&\lesssim \frac{|\tilde{Q}|^{1/2}}{|\tilde{e}_6|^{1/2}} \frac{|\tilde{e}_6|^{1/2}}{|\tilde{Q}|^{1/2}}(\|\delta\boldsymbol{\tilde{v}}_1\|_{L^2(\tilde{Q})} + h_{\tilde{Q}}|\delta\boldsymbol{\tilde{v}}_1|_{H^1(\tilde{Q})} )\\
		&\lesssim |\boldsymbol{\tilde{v}}_1|_{H^1(\tilde{Q})}.
	\end{align*}
	Therefore,
	\begin{equation}\label{eq:pi_m^1Q}
		|\tilde{\Pi}^{(1)}_{\tilde{M}} \boldsymbol{\tilde{v}}_1|_{H^1(\tilde{Q})} \lesssim |\boldsymbol{\tilde{v}}_1|_{H^1(\tilde{Q})}.
	\end{equation}
	and furthermore
	\begin{equation*}
		|\Pi^{(1)}_M \boldsymbol{v}|_{H^1(Q)} \lesssim |\tilde{\Pi}^{(1)}_{\tilde{M}} \boldsymbol{\tilde{v}}|_{H^1(\tilde{Q})}
		\lesssim |\boldsymbol{\tilde{v}}_1|_{H^1(\tilde{Q})}\lesssim |\bm{v}_1|_{H^1(Q)}.
	\end{equation*}
	Combining \eqref{eq:pi_m^1err} and \eqref{eq:pi_m^1Q}:
	\begin{equation}\label{eq:pi_m^1QT}
		|\Pi^{(1)}_M \boldsymbol{v}_1|_{H^1(T \cup Q)} \lesssim |\boldsymbol{v}_1|_{H^1(M)}.
	\end{equation}
	The proof completes by Fortin trick.
\end{proof}

\end{document}
